\begin{claim}{F.12}{Proved}{The common involution quotient and its winding equations}
Use the actual action and source in F.11.  Put
\[
 \tau=(1\ 5)(3\ 4),\qquad
 \beta=(1\ 3)(4\ 5),\qquad
 Y=\mathbf P^1\setminus\{0,\infty,\pm\mu_5\},
               \qquad t=s^2 .
                                                               \tag{F.12.1}
\]
The map $f$ in F.11 induces a split injection on
abelianizations, including their full profinite
completions.  Consequently its completed kernel
lies in the closed commutator subgroup of the
specified $\widehat F_6$.

There is an explicit, fully injective map
\[
 \widehat\phi:\widehat{\pi_1(Y,s_0)}
                     =\widehat F_{11}
                   \longrightarrow\widehat{P\Gamma}_{0,8},
 \qquad
 s_0=i(\theta^2-2)\zeta,\quad s_0^2=t_0,
                                                               \tag{F.12.2}
\]
induced by the configuration
\[
 \phi(s)=(0,\infty,1,v_1(s),\ldots,v_5(s)),
 \qquad
 v_{j+1}(s)=
       \left(\frac{\zeta^j-s}{\zeta^j+s}\right)^2
                         \quad(0\leq j\leq4).
                                                               \tag{F.12.3}
\]
The actual quotient action of $\beta$ induces
the following based involution on this normalized
eight-point moduli space:
\[
 R(v_1,v_2,v_3,v_4,v_5)
       =(v_2/v_3,\ v_2,\ v_2/v_1,\ v_2/v_5,\ v_2/v_4).
                                                               \tag{F.12.4}
\]
The residual centralizer in F.7 is exactly the
image, under the specified covering inclusion,
of the full equalizer
\[
 \mathcal E=
 \{u\in\widehat{\pi_1(Y,s_0)}:
                    \widehat\phi(u)=R_*\widehat\phi(u)\}.
                                                               \tag{F.12.5}
\]
In particular this identifies all the groups and
maps in the remaining fixed-intersection problem.

Write $e_j^\pm$ for the positively oriented source
meridian classes at $\pm\zeta^j$.  The abelian
image of every $u\in\mathcal E$ has the form
\[
        [u]=\sum_{j=0}^4 a_j(e_j^+-e_j^-),\qquad
        a_0-a_1+a_2=0,\quad -a_1+a_3+a_4=0
                  \quad\hbox{in }\widehat\Z.
                                                               \tag{F.12.6}
\]
These are exactly the restrictions of the twenty
target winding equalities to that anti-invariant
source submodule.  Their kernel is the explicitly
specified free rank-three $\widehat\Z$-module
in the proof; no equality with the actual image
of $\mathcal E$ is asserted.

For the element $w$ in F.11, the lift to this
double-cover source has
$[w]=e_1^- -e_1^+$.
Its first target winding defect is exactly four.
Hence its entire procyclic closure meets
$\mathcal E$ trivially, by this specific
abelian quotient computation.
Triviality of the full equalizer (F.12.5)
remains unresolved.
\end{claim}

\begin{proof}
\emph{Six exact characters for the seven-point map.}
On the target of F.11.2, use the last labelled
point as infinity and normalize the first two
finite points to $0,1$.  Denote the resulting
six finite coordinates by $z_0,\ldots,z_5$,
so $z_0=0,z_1=1$.  Their pullbacks under $f$ are
\[
 (z_0,\ldots,z_5)=
 \frac{(1,\zeta^2,\zeta^4,t,t\zeta,t\zeta^4)-1}
                         {\zeta^2-1}.
                                                               \tag{F.12.7}
\]
For distinct $i,j$, the function $z_i-z_j$
is a globally defined nonzero function on
this labelled moduli space, so its winding
defines a homomorphism to $\Z$.
Take the six pairs, in this order,
\[
            (3,4),\ (3,0),\ (5,0),\
                         (3,1),\ (5,1),\ (4,0).
                                                               \tag{F.12.8}
\]
The six pullbacks are nonzero constant multiples
of
\[
             t,\quad t-1,\quad t-\zeta,\quad
                       t-\zeta^2,\quad t-\zeta^3,\quad t-\zeta^4.
                                                               \tag{F.12.9}
\]
Thus the induced homomorphism
$\ell:P\Gamma_{0,7}\to\Z^6$ satisfies
$\ell f_*=\operatorname{ab}$ on
$\pi_1(X_5)$, in its ordinary finite-puncture
meridian basis.  Completion gives the same
identity with $\widehat\Z^6$.
This proves the asserted split injection and
the closed-commutator containment.
In particular the residual centralizer, which
belongs to this kernel by F.7.9 and F.11,
lies in the index-two subgroup corresponding
to $t=s^2$.
The source is exactly $Y$ in (F.12.1);
finite-index subgroups have their full
induced profinite topology.

\emph{The actual common-involution quotient.}
For the family in F.11, the involution
$\tau_t(z)=t/z$ has fixed points $s,-s$.
On $Y$, label them in that order.
The actual quotient map, including these
branch labels, is
\[
                  Q_s(z)=\left(\frac{z-s}{z+s}\right)^2 .
                                                               \tag{F.12.10}
\]
It sends the two branch points to $0,\infty$.
The original pole pair $\{0,\infty\}$ goes to
$1$.  Its other five fibres on the original
twelve punctures are
\[
    \{\zeta^j,s^2\zeta^{-j}\},\qquad 0\leq j\leq4,
\]
and their images are precisely (F.12.3).
No two of the eight target points coincide
for $s\in Y$.
Indeed $v_{j+1}$ is $0$ or infinity only at
$s=\zeta^j$ or $-\zeta^j$, and equals $1$
only at $s=0,\infty$.
For $i\ne j$, the other collision equation
is $s^2=\zeta^{i+j}$, whose two solutions
again lie in $\pm\mu_5$.

Here is the full completion comparison for
this quotient.  The discrete normalizer
sequence for $\langle\tau\rangle$ is the
specified sequence F.7.5:
\[
 1\longrightarrow\langle\tau\rangle
 \longrightarrow N_M(\langle\tau\rangle)
 \longrightarrow Q\longrightarrow1 .
                                                               \tag{F.12.11}
\]
Here $Q$ is the group of the quotient sphere
with its six punctures and two branch points,
with their types retained.  It has finite
index in the mapping class group with these
decorations.  The prime-order comparison and
the full topology in F.7.4--5 identify its
completed source with the actual completed
normalizer.  Consequently completion of
(F.12.11) has exactly the displayed finite
deck kernel and has target $\widehat Q$.
The deck involution permutes the twelve
original punctures.  Its intersection with
the completed pure centralizer is therefore
trivial, as detected by the finite permutation
quotient.  Descent is injective on
$C_{\widehat P}(\tau)$.
The finite-index inclusion of $Q$ supplies
its full embedding in the decorated
eight-point mapping class completion.

The map from $\widehat{\pi_1(Y)}$ into
$\widehat P$ is the index-two covering
inclusion followed by F.11.10.
It is injective and its image centralizes
$\tau$.  On $Y$ the two branch labels do
not exchange, so the quotient map lands
in $\widehat{P\Gamma}_{0,8}$.
Its composite is exactly
$\widehat\phi$ in (F.12.2), by (F.12.10).
Both components are injective.  This
proves the full injectivity in (F.12.2),
with all branch labels retained.

\emph{The exact second involution and based naturality.}
In the twelve-label table of F.11, the six
$\tau$-orbits, in the order used for the
quotient punctures $1,v_1,\ldots,v_5$, are
\[
 \{0,10\},\quad\{1,4\},\quad\{3,8\},\quad
              \{5,11\},\quad\{6,9\},\quad\{2,7\}.
                                                               \tag{F.12.12}
\]
The element $\beta=(1\ 3)(4\ 5)$ commutes
with $\tau$, is outside $D_{10}$, and
together with $D_{10}$ generates $A_5$.
Its permutation on these six orbits is
\[
              p=(0\ 2)(1\ 3)(4\ 5).
                                                               \tag{F.12.13}
\]
It exchanges the two branch points.
These facts are checked in the same actual
M\"obius action as F.11, not merely in a
separate permutation representation.
Specifically the word
$b^{-1}aba b^{-1}ab$ has the named action
of $\beta$, sends $s_0$ to $-s_0$ and
$-s_0$ to $s_0$, and in the coordinate
$(z-s_0)/(z+s_0)$ has the form $c/z$.
The certificate verifies $c^2=v_2(s_0)$.
Thus its actual quotient representative is
$V\mapsto v_2(s_0)/V$.

For general normalized quotient coordinates
write $v_0=1$.
After applying the label permutation and
renormalizing the branch values to
$0,\infty$ and puncture zero to $1$,
the new coordinate at puncture $i$ is
$v_{p(0)}/v_{p(i)}$.
This is precisely (F.12.4).
At $s_0$ the actual quotient representative
just specified proves
$R(\phi(s_0))=\phi(s_0)$.
The based action $R_*$ is therefore fixed
without an unspecified base path, and
agrees with conjugation by that quotient
mapping class.
The source basepoint comparison with F.11
uses the unique lift, starting at $s_0$,
of the already specified straight segment
from $t_0$ to $2+2i$.

Let $\jmath:\widehat{\pi_1(Y)}\to\widehat P$
denote the actual covering inclusion above.
For the residual centralizer, (F.12.7)--(F.12.9)
show that its elements have unique preimages
under $\jmath$.
The descent of each such element commutes
with the quotient of $\beta$, proving that
its preimage lies in $\mathcal E$.
Conversely, for $u\in\mathcal E$, both
$\jmath(u)$ and
$\beta\jmath(u)\beta^{-1}$ lie in
$C_{\widehat P}(\tau)$.
Their descents agree by (F.12.5).
The injectivity proved from (F.12.11)
therefore gives
$\beta\jmath(u)\beta^{-1}=\jmath(u)$.
The element $\jmath(u)$ already centralizes
$D_{10}$, so it centralizes their generated
group $A_5$.
This proves both directions of (F.12.5);
the resulting isomorphism is continuous
with continuous inverse on these compact
subgroups.

\emph{All winding equations on the source homology.}
Use the finite source meridians at
$0,\zeta^j,-\zeta^j$ as a basis of
$H_1(Y;\widehat\Z)$, omitting the infinity
meridian.  The map $s\mapsto s^2$ sends
the zero meridian to twice the zero
meridian downstairs, and sends each
$e_j^\pm$ to the downstairs meridian at
$\zeta^{2j}$.
The downstairs abelianization of an element
of $\mathcal E$ vanishes by the first part
of the proof.  Since multiplication by
two on $\widehat\Z$ is injective, its
zero coefficient is zero and its two
coefficients above each $\zeta^{2j}$ sum
to zero.  This gives the first expression
in (F.12.6).

For clarity the remaining calculation is
entirely integral.  Take the twenty winding
characters on the normalized eight-point
space:
\[
 L_i=\operatorname{wind}(v_i),\quad
 M_i=\operatorname{wind}(v_i-1)\quad(1\leq i\leq5),
 \qquad
 N_{ij}=\operatorname{wind}(v_i-v_j)\quad(1\leq i<j\leq5).
                                                               \tag{F.12.14}
\]
These give its usual integral first-cohomology
basis.  One can see the basis directly by
forgetting the last variable in turn: the
fibre is the plane with the earlier points
and $0,1$ removed, its meridians are detected
by the newly added differences, and pure
monodromy fixes their abelian classes.
The successive ranks are $2,3,4,5,6$.

The exact functions in (F.12.3) satisfy
\[
 \begin{split}
 v_{j+1}-1
     &=\frac{-4\zeta^j s}{(\zeta^j+s)^2},\\
 v_{i+1}-v_{j+1}
     &=\frac{4s(\zeta^i-\zeta^j)(\zeta^{i+j}-s^2)}
             {(\zeta^i+s)^2(\zeta^j+s)^2}
                         \qquad(i\ne j).
 \end{split}                                                   \tag{F.12.15}
\]
Their zero and pole orders at the twelve
source punctures determine the pullbacks.
On an anti-invariant class
$\sum a_j(e_j^+-e_j^-)$ the resulting values
are
\[
       L_{j+1}=4a_j,\qquad M_{j+1}=2a_j,\qquad
                   N_{i+1,j+1}=2(a_i+a_j).
                                                               \tag{F.12.16}
\]
The extra zeros of the last expression,
at the two points with $s^2=\zeta^{i+j}$,
cancel on an anti-invariant class.

Put $A_0=0$ and $A_{j+1}=a_j$.
For $v_0=1$, the winding of any difference
$v_i-v_j$ is uniformly $2(A_i+A_j)$,
and that of $v_i$ is $4A_i$.
Under $R$, these functions are transformed
using $v_i\mapsto v_2/v_{p(i)}$.
The twenty equalities of winding values
are consequently equivalent to
\[
                   A_i+A_{p(i)}=A_2
                             \qquad(0\leq i\leq5).            \tag{F.12.17}
\]
Here necessity already follows from
$M_1$ and $M_4$, whose defects are
$2(a_0-a_1+a_2)$ and
$2(-a_1+a_3+a_4)$.
Multiplication by two is injective.
For sufficiency, substitute (F.12.17)
in the winding of $v_2/v_{p(i)}$,
$(v_2-v_{p(i)})/v_{p(i)}$, and
$v_2(v_{p(j)}-v_{p(i)})/(v_{p(i)}v_{p(j)})$;
these are exactly the three types in
(F.12.14).
The only two nontrivial relations in
(F.12.17) are (F.12.6).
Their kernel has the following integral
basis, which also gives its full
$\widehat\Z$-basis:
\[
          (-1,0,1,0,0),\qquad
          (1,1,0,1,0),\qquad
          (1,1,0,0,1).
                                                               \tag{F.12.18}
\]
These equations are necessary for
membership in $\mathcal E$; their kernel
is only the module defined by the
displayed winding equations.

\emph{The concrete detector for the element of F.11.}
At $t_0$ the chosen root is
$s_0=i(\theta^2-2)\zeta$.
The lift of the straight path to $b=2+2i$
ends at the negative square root of $b$:
throughout that path the real part of
$t$ is positive, and $s_0$ is the root
with negative real part.
The whisker and the small square for
$m_{\zeta^2}$ lie in the upper half-plane.
Continuation on this sheet therefore
lifts that meridian to a positive
meridian at $-\zeta$.
Its lift on the other sheet is a meridian
at $+\zeta$.

The second commutator factor
$m_\zeta m_{\zeta^3}m_0$ has winding one
about zero, and hence exchanges the
two sheets.  In the lifted commutator,
its path and its reverse cancel in
abelianization.  The two remaining
meridians give exactly
\[
                         [w]=e_1^- -e_1^+ .
                                                               \tag{F.12.19}
\]
By (F.12.4), $R^*L_1=L_2-L_3$.
The winding difference for this class
is therefore
\[
       (L_1-R^*L_1)\widehat\phi(w)
               =4(a_0-a_1+a_2)=4.
                                                               \tag{F.12.20}
\]
For a full profinite power $w^u$, the
continuous extension of this character
has value $4u\in\widehat\Z$.
It vanishes only for $u=0$.
This is an explicit family of finite
cyclic detectors for the procyclic
exclusion, and uses the actual common
involution quotient.

The executed certificate
\path{code/part_a/c5_common_reflection_f12.py}
verifies the common finite labels, the
M\"obius involution and its branch
normalization, the six-character splitting,
the twenty-by-twelve integral divisor
matrix, the exact involution on the
twenty winding characters, and the
restriction to the two rows in (F.12.6).
It also records the winding defect four.
Its larger unreduced matrices have source
rank ten and fixed-difference rank six;
no rank calculation is used to infer
triviality of the full equalizer.
\end{proof}
